\uuid{nNt2}
\exo7id{4607}
\titre{exo7 4607}
\auteur{quercia}
\organisation{exo7}
\datecreate{2010-03-14}
\isIndication{false}
\isCorrection{true}
\chapitre{Série entière}
\sousChapitre{Intégrales}
\cpge{mpsi-mp2i.procedes-sommatoires-discrets.a, pcsi.series-numeriques.a, ptsi.series-numeriques.a, mpsi-mp2i.integration.c, pcsi.integration.b, ptsi.integration.b, bcpst1.integration-segment.a, mp-mpi.series-entieres.d, pc.series-entieres.c, psi.series-entieres.c, pt.series-entieres.c, mp-mpi.series-entieres.a, pc.series-entieres.a, psi.series-entieres.a, pt.series-entieres.a, bcpst2.series-reelles}
\module{Analyse}
\niveau{L2}
\difficulte{}

\contenu{
\begin{enumerate}
    \item \question{A l'aide d'un développement en série entière, montrer que
    $ \int_{t=0}^1 t^t\,d t = \sum_{n=1}^\infty \frac{(-1)^{n-1}}{n^n}$.}
\reponse{$t^t = \exp(t\ln t) = \sum_{k=0}^\infty \frac{t^k\ln^k t}{k!}$.}
    \item \question{Calculer la valeur commune des deux membres à $10^{-5}$ près.}
\reponse{0.78343}
\end{enumerate}
}
